% RQ2 — Results subsection (prose + summary box)
% Requires: \usepackage{tcolorbox}
% Usage: % RQ2 — Results subsection (prose + summary box)
% Requires: \usepackage{tcolorbox}
% Usage: % RQ2 — Results subsection (prose + summary box)
% Requires: \usepackage{tcolorbox}
% Usage: % RQ2 — Results subsection (prose + summary box)
% Requires: \usepackage{tcolorbox}
% Usage: \input{paper_tables/rq2_results_text.tex}
%         \input{paper_tables/rq2_results.tex}

\subsection{RQ2: RetroBug vs.\ Zero-Shot LLM Baselines}
\label{sec:results:rq2}
% --------------------------------------------------------------

Table~\ref{tab:rq2} compares RetroBug with three zero-shot LLM baselines
on fold~0 ($n=160$).
All baselines receive the same inputs---a masked commit message and the
full git diff---and produce a single-pass binary label without knowledge-base
retrieval, backward reasoning, or tool use.
Claude Sonnet~4.5 and GPT-4o follow the conventional commit-classification
prompt of Zeng~\textit{et al.}~\cite{zeng2025conventional}; CodeLlama-13b
uses an equivalent JSON-structured prompt with prose fallback parsing.

RetroBug achieves Bug-Fix F1~=~0.776 and MCC~=~0.698, outperforming every
zero-shot baseline on all five reported metrics.
The strongest zero-shot competitor is GPT-4o (F1~=~0.629, MCC~=~0.493),
followed by Claude Sonnet~4.5 (F1~=~0.598, MCC~=~0.446) and CodeLlama-13b
(F1~=~0.431, MCC~=~0.171).
Relative to GPT-4o, RetroBug improves Bug-Fix F1 by 14.7 points and MCC
by 0.205 points; relative to Claude, the gains are 17.8 and 0.252 points
respectively.

The performance gap is driven primarily by \emph{precision}.
All three zero-shot models exhibit a Bug-Fix bias: they achieve high recall
(0.625--0.800) but low precision (0.329--0.571), misclassifying a large
share of non-bug commits as bug-fixes.
Claude Sonnet~4.5, for example, labels 35 of 120 non-bug commits as
Bug-Fix (precision 0.478), while still recovering 32 of 40 true bug-fix
commits (recall 0.800).
CodeLlama-13b amplifies this pattern most severely: it predicts Bug-Fix
for 76 commits but only 25 are correct (precision 0.329), yielding accuracy
of just 0.585 despite moderate recall (0.625).
GPT-4o balances the trade-off best among zero-shot models (precision
0.571, recall 0.700), but still falls 16.2 precision points below
RetroBug (0.733).

RetroBug's advantage over single-pass prompting indicates that multi-step
backward CBR, knowledge-base retrieval, and structured diff analysis add
value beyond exposing the same evidence to a capable LLM in one shot.
The agent recovers 33 of 40 bug-fix commits (recall 0.825) while limiting
false positives to 12 of 120 non-bug commits---a precision--recall profile
that none of the zero-shot baselines match.
We evaluate LLM baselines on fold~0 only due to API cost and inference
latency; RetroBug results on the same fold are directly comparable because
they use the identical 160-commit test split.

\input{paper_tables/rq2_results.tex}

\begin{tcolorbox}[colback=gray!8, colframe=gray!50,
  boxrule=0.5pt, arc=2pt, left=4pt, right=4pt,
  top=3pt, bottom=3pt]
\textbf{Answer to RQ2:}
On fold~0, RetroBug (Bug-Fix F1~=~0.776, MCC~=~0.698) outperforms the
best zero-shot LLM baseline, GPT-4o (F1~=~0.629, MCC~=~0.493), by 14.7
F1 points and 0.205 MCC points.
Zero-shot models achieve competitive recall but suffer from low Bug-Fix
precision (0.329--0.571), whereas RetroBug combines high recall (0.825)
with substantially higher precision (0.733).
These results show that agentic backward CBR and tool use improve
classification beyond single-pass LLM prompting on the same inputs.
\end{tcolorbox}

%         % RQ2 — RetroBug vs CCS baselines (fold 0, per-class metrics)
% Requires: \usepackage{booktabs}
% Usage: \input{paper_tables/rq2_results.tex}

\begin{table}[t]
\centering
\scriptsize
\caption{RQ2: Per-class performance of RetroBug versus CCS zero-shot
LLM baselines on fold~0 ($n=160$).
NB~=~Non-Bug-Fix, BF~=~Bug-Fix.}
\label{tab:rq2}
\setlength{\tabcolsep}{3pt}
\renewcommand{\arraystretch}{0.9}
\begin{tabular}{@{}l ccc ccc cc@{}}
\toprule
 & \multicolumn{3}{c}{\textbf{Non-Bug-Fix (NB)}} & \multicolumn{3}{c}{\textbf{Bug-Fix (BF)}} & & \\
\cmidrule(lr){2-4} \cmidrule(lr){5-7}
\textbf{Method} & \textbf{P} & \textbf{R} & \textbf{F1} & \textbf{P} & \textbf{R} & \textbf{F1} & \textbf{Acc.} & \textbf{MCC} \\
\midrule
CCS (Claude)~\cite{zeng2025conventional}
  & 0.914 & 0.708 & 0.798 & 0.478 & 0.800 & 0.598 & 0.731 & 0.446 \\
\midrule
CCS (GPT-4o)~\cite{zeng2025conventional}
  & 0.892 & 0.825 & 0.857 & 0.571 & 0.700 & 0.629 & 0.794 & 0.493 \\
\midrule
CCS (CodeLlama-13b)~\cite{zeng2025conventional}
  & 0.819 & 0.571 & 0.673 & 0.329 & 0.625 & 0.431 & 0.585 & 0.171 \\
\midrule
\textbf{RetroBug (ours)}
  & \textbf{0.939} & \textbf{0.900} & \textbf{0.919} & \textbf{0.733} & \textbf{0.825} & \textbf{0.776} & \textbf{0.881} & \textbf{0.698} \\
\bottomrule
\end{tabular}

\vspace{0.3em}
{\footnotesize CodeLlama-13b: $n=159$ (one unparseable response excluded).}
\end{table}


\subsection{RQ2: RetroBug vs.\ Zero-Shot LLM Baselines}
\label{sec:results:rq2}
% --------------------------------------------------------------

Table~\ref{tab:rq2} compares RetroBug with three zero-shot LLM baselines
on fold~0 ($n=160$).
All baselines receive the same inputs---a masked commit message and the
full git diff---and produce a single-pass binary label without knowledge-base
retrieval, backward reasoning, or tool use.
Claude Sonnet~4.5 and GPT-4o follow the conventional commit-classification
prompt of Zeng~\textit{et al.}~\cite{zeng2025conventional}; CodeLlama-13b
uses an equivalent JSON-structured prompt with prose fallback parsing.

RetroBug achieves Bug-Fix F1~=~0.776 and MCC~=~0.698, outperforming every
zero-shot baseline on all five reported metrics.
The strongest zero-shot competitor is GPT-4o (F1~=~0.629, MCC~=~0.493),
followed by Claude Sonnet~4.5 (F1~=~0.598, MCC~=~0.446) and CodeLlama-13b
(F1~=~0.431, MCC~=~0.171).
Relative to GPT-4o, RetroBug improves Bug-Fix F1 by 14.7 points and MCC
by 0.205 points; relative to Claude, the gains are 17.8 and 0.252 points
respectively.

The performance gap is driven primarily by \emph{precision}.
All three zero-shot models exhibit a Bug-Fix bias: they achieve high recall
(0.625--0.800) but low precision (0.329--0.571), misclassifying a large
share of non-bug commits as bug-fixes.
Claude Sonnet~4.5, for example, labels 35 of 120 non-bug commits as
Bug-Fix (precision 0.478), while still recovering 32 of 40 true bug-fix
commits (recall 0.800).
CodeLlama-13b amplifies this pattern most severely: it predicts Bug-Fix
for 76 commits but only 25 are correct (precision 0.329), yielding accuracy
of just 0.585 despite moderate recall (0.625).
GPT-4o balances the trade-off best among zero-shot models (precision
0.571, recall 0.700), but still falls 16.2 precision points below
RetroBug (0.733).

RetroBug's advantage over single-pass prompting indicates that multi-step
backward CBR, knowledge-base retrieval, and structured diff analysis add
value beyond exposing the same evidence to a capable LLM in one shot.
The agent recovers 33 of 40 bug-fix commits (recall 0.825) while limiting
false positives to 12 of 120 non-bug commits---a precision--recall profile
that none of the zero-shot baselines match.
We evaluate LLM baselines on fold~0 only due to API cost and inference
latency; RetroBug results on the same fold are directly comparable because
they use the identical 160-commit test split.

% RQ2 — RetroBug vs CCS baselines (fold 0, per-class metrics)
% Requires: \usepackage{booktabs}
% Usage: \input{paper_tables/rq2_results.tex}

\begin{table}[t]
\centering
\scriptsize
\caption{RQ2: Per-class performance of RetroBug versus CCS zero-shot
LLM baselines on fold~0 ($n=160$).
NB~=~Non-Bug-Fix, BF~=~Bug-Fix.}
\label{tab:rq2}
\setlength{\tabcolsep}{3pt}
\renewcommand{\arraystretch}{0.9}
\begin{tabular}{@{}l ccc ccc cc@{}}
\toprule
 & \multicolumn{3}{c}{\textbf{Non-Bug-Fix (NB)}} & \multicolumn{3}{c}{\textbf{Bug-Fix (BF)}} & & \\
\cmidrule(lr){2-4} \cmidrule(lr){5-7}
\textbf{Method} & \textbf{P} & \textbf{R} & \textbf{F1} & \textbf{P} & \textbf{R} & \textbf{F1} & \textbf{Acc.} & \textbf{MCC} \\
\midrule
CCS (Claude)~\cite{zeng2025conventional}
  & 0.914 & 0.708 & 0.798 & 0.478 & 0.800 & 0.598 & 0.731 & 0.446 \\
\midrule
CCS (GPT-4o)~\cite{zeng2025conventional}
  & 0.892 & 0.825 & 0.857 & 0.571 & 0.700 & 0.629 & 0.794 & 0.493 \\
\midrule
CCS (CodeLlama-13b)~\cite{zeng2025conventional}
  & 0.819 & 0.571 & 0.673 & 0.329 & 0.625 & 0.431 & 0.585 & 0.171 \\
\midrule
\textbf{RetroBug (ours)}
  & \textbf{0.939} & \textbf{0.900} & \textbf{0.919} & \textbf{0.733} & \textbf{0.825} & \textbf{0.776} & \textbf{0.881} & \textbf{0.698} \\
\bottomrule
\end{tabular}

\vspace{0.3em}
{\footnotesize CodeLlama-13b: $n=159$ (one unparseable response excluded).}
\end{table}


\begin{tcolorbox}[colback=gray!8, colframe=gray!50,
  boxrule=0.5pt, arc=2pt, left=4pt, right=4pt,
  top=3pt, bottom=3pt]
\textbf{Answer to RQ2:}
On fold~0, RetroBug (Bug-Fix F1~=~0.776, MCC~=~0.698) outperforms the
best zero-shot LLM baseline, GPT-4o (F1~=~0.629, MCC~=~0.493), by 14.7
F1 points and 0.205 MCC points.
Zero-shot models achieve competitive recall but suffer from low Bug-Fix
precision (0.329--0.571), whereas RetroBug combines high recall (0.825)
with substantially higher precision (0.733).
These results show that agentic backward CBR and tool use improve
classification beyond single-pass LLM prompting on the same inputs.
\end{tcolorbox}

%         % RQ2 — RetroBug vs CCS baselines (fold 0, per-class metrics)
% Requires: \usepackage{booktabs}
% Usage: % RQ2 — RetroBug vs CCS baselines (fold 0, per-class metrics)
% Requires: \usepackage{booktabs}
% Usage: \input{paper_tables/rq2_results.tex}

\begin{table}[t]
\centering
\scriptsize
\caption{RQ2: Per-class performance of RetroBug versus CCS zero-shot
LLM baselines on fold~0 ($n=160$).
NB~=~Non-Bug-Fix, BF~=~Bug-Fix.}
\label{tab:rq2}
\setlength{\tabcolsep}{3pt}
\renewcommand{\arraystretch}{0.9}
\begin{tabular}{@{}l ccc ccc cc@{}}
\toprule
 & \multicolumn{3}{c}{\textbf{Non-Bug-Fix (NB)}} & \multicolumn{3}{c}{\textbf{Bug-Fix (BF)}} & & \\
\cmidrule(lr){2-4} \cmidrule(lr){5-7}
\textbf{Method} & \textbf{P} & \textbf{R} & \textbf{F1} & \textbf{P} & \textbf{R} & \textbf{F1} & \textbf{Acc.} & \textbf{MCC} \\
\midrule
CCS (Claude)~\cite{zeng2025conventional}
  & 0.914 & 0.708 & 0.798 & 0.478 & 0.800 & 0.598 & 0.731 & 0.446 \\
\midrule
CCS (GPT-4o)~\cite{zeng2025conventional}
  & 0.892 & 0.825 & 0.857 & 0.571 & 0.700 & 0.629 & 0.794 & 0.493 \\
\midrule
CCS (CodeLlama-13b)~\cite{zeng2025conventional}
  & 0.819 & 0.571 & 0.673 & 0.329 & 0.625 & 0.431 & 0.585 & 0.171 \\
\midrule
\textbf{RetroBug (ours)}
  & \textbf{0.939} & \textbf{0.900} & \textbf{0.919} & \textbf{0.733} & \textbf{0.825} & \textbf{0.776} & \textbf{0.881} & \textbf{0.698} \\
\bottomrule
\end{tabular}

\vspace{0.3em}
{\footnotesize CodeLlama-13b: $n=159$ (one unparseable response excluded).}
\end{table}


\begin{table}[t]
\centering
\scriptsize
\caption{RQ2: Per-class performance of RetroBug versus CCS zero-shot
LLM baselines on fold~0 ($n=160$).
NB~=~Non-Bug-Fix, BF~=~Bug-Fix.}
\label{tab:rq2}
\setlength{\tabcolsep}{3pt}
\renewcommand{\arraystretch}{0.9}
\begin{tabular}{@{}l ccc ccc cc@{}}
\toprule
 & \multicolumn{3}{c}{\textbf{Non-Bug-Fix (NB)}} & \multicolumn{3}{c}{\textbf{Bug-Fix (BF)}} & & \\
\cmidrule(lr){2-4} \cmidrule(lr){5-7}
\textbf{Method} & \textbf{P} & \textbf{R} & \textbf{F1} & \textbf{P} & \textbf{R} & \textbf{F1} & \textbf{Acc.} & \textbf{MCC} \\
\midrule
CCS (Claude)~\cite{zeng2025conventional}
  & 0.914 & 0.708 & 0.798 & 0.478 & 0.800 & 0.598 & 0.731 & 0.446 \\
\midrule
CCS (GPT-4o)~\cite{zeng2025conventional}
  & 0.892 & 0.825 & 0.857 & 0.571 & 0.700 & 0.629 & 0.794 & 0.493 \\
\midrule
CCS (CodeLlama-13b)~\cite{zeng2025conventional}
  & 0.819 & 0.571 & 0.673 & 0.329 & 0.625 & 0.431 & 0.585 & 0.171 \\
\midrule
\textbf{RetroBug (ours)}
  & \textbf{0.939} & \textbf{0.900} & \textbf{0.919} & \textbf{0.733} & \textbf{0.825} & \textbf{0.776} & \textbf{0.881} & \textbf{0.698} \\
\bottomrule
\end{tabular}

\vspace{0.3em}
{\footnotesize CodeLlama-13b: $n=159$ (one unparseable response excluded).}
\end{table}


\subsection{RQ2: RetroBug vs.\ Zero-Shot LLM Baselines}
\label{sec:results:rq2}
% --------------------------------------------------------------

Table~\ref{tab:rq2} compares RetroBug with three zero-shot LLM baselines
on fold~0 ($n=160$).
All baselines receive the same inputs---a masked commit message and the
full git diff---and produce a single-pass binary label without knowledge-base
retrieval, backward reasoning, or tool use.
Claude Sonnet~4.5 and GPT-4o follow the conventional commit-classification
prompt of Zeng~\textit{et al.}~\cite{zeng2025conventional}; CodeLlama-13b
uses an equivalent JSON-structured prompt with prose fallback parsing.

RetroBug achieves Bug-Fix F1~=~0.776 and MCC~=~0.698, outperforming every
zero-shot baseline on all five reported metrics.
The strongest zero-shot competitor is GPT-4o (F1~=~0.629, MCC~=~0.493),
followed by Claude Sonnet~4.5 (F1~=~0.598, MCC~=~0.446) and CodeLlama-13b
(F1~=~0.431, MCC~=~0.171).
Relative to GPT-4o, RetroBug improves Bug-Fix F1 by 14.7 points and MCC
by 0.205 points; relative to Claude, the gains are 17.8 and 0.252 points
respectively.

The performance gap is driven primarily by \emph{precision}.
All three zero-shot models exhibit a Bug-Fix bias: they achieve high recall
(0.625--0.800) but low precision (0.329--0.571), misclassifying a large
share of non-bug commits as bug-fixes.
Claude Sonnet~4.5, for example, labels 35 of 120 non-bug commits as
Bug-Fix (precision 0.478), while still recovering 32 of 40 true bug-fix
commits (recall 0.800).
CodeLlama-13b amplifies this pattern most severely: it predicts Bug-Fix
for 76 commits but only 25 are correct (precision 0.329), yielding accuracy
of just 0.585 despite moderate recall (0.625).
GPT-4o balances the trade-off best among zero-shot models (precision
0.571, recall 0.700), but still falls 16.2 precision points below
RetroBug (0.733).

RetroBug's advantage over single-pass prompting indicates that multi-step
backward CBR, knowledge-base retrieval, and structured diff analysis add
value beyond exposing the same evidence to a capable LLM in one shot.
The agent recovers 33 of 40 bug-fix commits (recall 0.825) while limiting
false positives to 12 of 120 non-bug commits---a precision--recall profile
that none of the zero-shot baselines match.
We evaluate LLM baselines on fold~0 only due to API cost and inference
latency; RetroBug results on the same fold are directly comparable because
they use the identical 160-commit test split.

% RQ2 — RetroBug vs CCS baselines (fold 0, per-class metrics)
% Requires: \usepackage{booktabs}
% Usage: % RQ2 — RetroBug vs CCS baselines (fold 0, per-class metrics)
% Requires: \usepackage{booktabs}
% Usage: \input{paper_tables/rq2_results.tex}

\begin{table}[t]
\centering
\scriptsize
\caption{RQ2: Per-class performance of RetroBug versus CCS zero-shot
LLM baselines on fold~0 ($n=160$).
NB~=~Non-Bug-Fix, BF~=~Bug-Fix.}
\label{tab:rq2}
\setlength{\tabcolsep}{3pt}
\renewcommand{\arraystretch}{0.9}
\begin{tabular}{@{}l ccc ccc cc@{}}
\toprule
 & \multicolumn{3}{c}{\textbf{Non-Bug-Fix (NB)}} & \multicolumn{3}{c}{\textbf{Bug-Fix (BF)}} & & \\
\cmidrule(lr){2-4} \cmidrule(lr){5-7}
\textbf{Method} & \textbf{P} & \textbf{R} & \textbf{F1} & \textbf{P} & \textbf{R} & \textbf{F1} & \textbf{Acc.} & \textbf{MCC} \\
\midrule
CCS (Claude)~\cite{zeng2025conventional}
  & 0.914 & 0.708 & 0.798 & 0.478 & 0.800 & 0.598 & 0.731 & 0.446 \\
\midrule
CCS (GPT-4o)~\cite{zeng2025conventional}
  & 0.892 & 0.825 & 0.857 & 0.571 & 0.700 & 0.629 & 0.794 & 0.493 \\
\midrule
CCS (CodeLlama-13b)~\cite{zeng2025conventional}
  & 0.819 & 0.571 & 0.673 & 0.329 & 0.625 & 0.431 & 0.585 & 0.171 \\
\midrule
\textbf{RetroBug (ours)}
  & \textbf{0.939} & \textbf{0.900} & \textbf{0.919} & \textbf{0.733} & \textbf{0.825} & \textbf{0.776} & \textbf{0.881} & \textbf{0.698} \\
\bottomrule
\end{tabular}

\vspace{0.3em}
{\footnotesize CodeLlama-13b: $n=159$ (one unparseable response excluded).}
\end{table}


\begin{table}[t]
\centering
\scriptsize
\caption{RQ2: Per-class performance of RetroBug versus CCS zero-shot
LLM baselines on fold~0 ($n=160$).
NB~=~Non-Bug-Fix, BF~=~Bug-Fix.}
\label{tab:rq2}
\setlength{\tabcolsep}{3pt}
\renewcommand{\arraystretch}{0.9}
\begin{tabular}{@{}l ccc ccc cc@{}}
\toprule
 & \multicolumn{3}{c}{\textbf{Non-Bug-Fix (NB)}} & \multicolumn{3}{c}{\textbf{Bug-Fix (BF)}} & & \\
\cmidrule(lr){2-4} \cmidrule(lr){5-7}
\textbf{Method} & \textbf{P} & \textbf{R} & \textbf{F1} & \textbf{P} & \textbf{R} & \textbf{F1} & \textbf{Acc.} & \textbf{MCC} \\
\midrule
CCS (Claude)~\cite{zeng2025conventional}
  & 0.914 & 0.708 & 0.798 & 0.478 & 0.800 & 0.598 & 0.731 & 0.446 \\
\midrule
CCS (GPT-4o)~\cite{zeng2025conventional}
  & 0.892 & 0.825 & 0.857 & 0.571 & 0.700 & 0.629 & 0.794 & 0.493 \\
\midrule
CCS (CodeLlama-13b)~\cite{zeng2025conventional}
  & 0.819 & 0.571 & 0.673 & 0.329 & 0.625 & 0.431 & 0.585 & 0.171 \\
\midrule
\textbf{RetroBug (ours)}
  & \textbf{0.939} & \textbf{0.900} & \textbf{0.919} & \textbf{0.733} & \textbf{0.825} & \textbf{0.776} & \textbf{0.881} & \textbf{0.698} \\
\bottomrule
\end{tabular}

\vspace{0.3em}
{\footnotesize CodeLlama-13b: $n=159$ (one unparseable response excluded).}
\end{table}


\begin{tcolorbox}[colback=gray!8, colframe=gray!50,
  boxrule=0.5pt, arc=2pt, left=4pt, right=4pt,
  top=3pt, bottom=3pt]
\textbf{Answer to RQ2:}
On fold~0, RetroBug (Bug-Fix F1~=~0.776, MCC~=~0.698) outperforms the
best zero-shot LLM baseline, GPT-4o (F1~=~0.629, MCC~=~0.493), by 14.7
F1 points and 0.205 MCC points.
Zero-shot models achieve competitive recall but suffer from low Bug-Fix
precision (0.329--0.571), whereas RetroBug combines high recall (0.825)
with substantially higher precision (0.733).
These results show that agentic backward CBR and tool use improve
classification beyond single-pass LLM prompting on the same inputs.
\end{tcolorbox}

%         % RQ2 — RetroBug vs CCS baselines (fold 0, per-class metrics)
% Requires: \usepackage{booktabs}
% Usage: % RQ2 — RetroBug vs CCS baselines (fold 0, per-class metrics)
% Requires: \usepackage{booktabs}
% Usage: % RQ2 — RetroBug vs CCS baselines (fold 0, per-class metrics)
% Requires: \usepackage{booktabs}
% Usage: \input{paper_tables/rq2_results.tex}

\begin{table}[t]
\centering
\scriptsize
\caption{RQ2: Per-class performance of RetroBug versus CCS zero-shot
LLM baselines on fold~0 ($n=160$).
NB~=~Non-Bug-Fix, BF~=~Bug-Fix.}
\label{tab:rq2}
\setlength{\tabcolsep}{3pt}
\renewcommand{\arraystretch}{0.9}
\begin{tabular}{@{}l ccc ccc cc@{}}
\toprule
 & \multicolumn{3}{c}{\textbf{Non-Bug-Fix (NB)}} & \multicolumn{3}{c}{\textbf{Bug-Fix (BF)}} & & \\
\cmidrule(lr){2-4} \cmidrule(lr){5-7}
\textbf{Method} & \textbf{P} & \textbf{R} & \textbf{F1} & \textbf{P} & \textbf{R} & \textbf{F1} & \textbf{Acc.} & \textbf{MCC} \\
\midrule
CCS (Claude)~\cite{zeng2025conventional}
  & 0.914 & 0.708 & 0.798 & 0.478 & 0.800 & 0.598 & 0.731 & 0.446 \\
\midrule
CCS (GPT-4o)~\cite{zeng2025conventional}
  & 0.892 & 0.825 & 0.857 & 0.571 & 0.700 & 0.629 & 0.794 & 0.493 \\
\midrule
CCS (CodeLlama-13b)~\cite{zeng2025conventional}
  & 0.819 & 0.571 & 0.673 & 0.329 & 0.625 & 0.431 & 0.585 & 0.171 \\
\midrule
\textbf{RetroBug (ours)}
  & \textbf{0.939} & \textbf{0.900} & \textbf{0.919} & \textbf{0.733} & \textbf{0.825} & \textbf{0.776} & \textbf{0.881} & \textbf{0.698} \\
\bottomrule
\end{tabular}

\vspace{0.3em}
{\footnotesize CodeLlama-13b: $n=159$ (one unparseable response excluded).}
\end{table}


\begin{table}[t]
\centering
\scriptsize
\caption{RQ2: Per-class performance of RetroBug versus CCS zero-shot
LLM baselines on fold~0 ($n=160$).
NB~=~Non-Bug-Fix, BF~=~Bug-Fix.}
\label{tab:rq2}
\setlength{\tabcolsep}{3pt}
\renewcommand{\arraystretch}{0.9}
\begin{tabular}{@{}l ccc ccc cc@{}}
\toprule
 & \multicolumn{3}{c}{\textbf{Non-Bug-Fix (NB)}} & \multicolumn{3}{c}{\textbf{Bug-Fix (BF)}} & & \\
\cmidrule(lr){2-4} \cmidrule(lr){5-7}
\textbf{Method} & \textbf{P} & \textbf{R} & \textbf{F1} & \textbf{P} & \textbf{R} & \textbf{F1} & \textbf{Acc.} & \textbf{MCC} \\
\midrule
CCS (Claude)~\cite{zeng2025conventional}
  & 0.914 & 0.708 & 0.798 & 0.478 & 0.800 & 0.598 & 0.731 & 0.446 \\
\midrule
CCS (GPT-4o)~\cite{zeng2025conventional}
  & 0.892 & 0.825 & 0.857 & 0.571 & 0.700 & 0.629 & 0.794 & 0.493 \\
\midrule
CCS (CodeLlama-13b)~\cite{zeng2025conventional}
  & 0.819 & 0.571 & 0.673 & 0.329 & 0.625 & 0.431 & 0.585 & 0.171 \\
\midrule
\textbf{RetroBug (ours)}
  & \textbf{0.939} & \textbf{0.900} & \textbf{0.919} & \textbf{0.733} & \textbf{0.825} & \textbf{0.776} & \textbf{0.881} & \textbf{0.698} \\
\bottomrule
\end{tabular}

\vspace{0.3em}
{\footnotesize CodeLlama-13b: $n=159$ (one unparseable response excluded).}
\end{table}


\begin{table}[t]
\centering
\scriptsize
\caption{RQ2: Per-class performance of RetroBug versus CCS zero-shot
LLM baselines on fold~0 ($n=160$).
NB~=~Non-Bug-Fix, BF~=~Bug-Fix.}
\label{tab:rq2}
\setlength{\tabcolsep}{3pt}
\renewcommand{\arraystretch}{0.9}
\begin{tabular}{@{}l ccc ccc cc@{}}
\toprule
 & \multicolumn{3}{c}{\textbf{Non-Bug-Fix (NB)}} & \multicolumn{3}{c}{\textbf{Bug-Fix (BF)}} & & \\
\cmidrule(lr){2-4} \cmidrule(lr){5-7}
\textbf{Method} & \textbf{P} & \textbf{R} & \textbf{F1} & \textbf{P} & \textbf{R} & \textbf{F1} & \textbf{Acc.} & \textbf{MCC} \\
\midrule
CCS (Claude)~\cite{zeng2025conventional}
  & 0.914 & 0.708 & 0.798 & 0.478 & 0.800 & 0.598 & 0.731 & 0.446 \\
\midrule
CCS (GPT-4o)~\cite{zeng2025conventional}
  & 0.892 & 0.825 & 0.857 & 0.571 & 0.700 & 0.629 & 0.794 & 0.493 \\
\midrule
CCS (CodeLlama-13b)~\cite{zeng2025conventional}
  & 0.819 & 0.571 & 0.673 & 0.329 & 0.625 & 0.431 & 0.585 & 0.171 \\
\midrule
\textbf{RetroBug (ours)}
  & \textbf{0.939} & \textbf{0.900} & \textbf{0.919} & \textbf{0.733} & \textbf{0.825} & \textbf{0.776} & \textbf{0.881} & \textbf{0.698} \\
\bottomrule
\end{tabular}

\vspace{0.3em}
{\footnotesize CodeLlama-13b: $n=159$ (one unparseable response excluded).}
\end{table}


\subsection{RQ2: RetroBug vs.\ Zero-Shot LLM Baselines}
\label{sec:results:rq2}
% --------------------------------------------------------------

Table~\ref{tab:rq2} compares RetroBug with three zero-shot LLM baselines
on fold~0 ($n=160$).
All baselines receive the same inputs---a masked commit message and the
full git diff---and produce a single-pass binary label without knowledge-base
retrieval, backward reasoning, or tool use.
Claude Sonnet~4.5 and GPT-4o follow the conventional commit-classification
prompt of Zeng~\textit{et al.}~\cite{zeng2025conventional}; CodeLlama-13b
uses an equivalent JSON-structured prompt with prose fallback parsing.

RetroBug achieves Bug-Fix F1~=~0.776 and MCC~=~0.698, outperforming every
zero-shot baseline on all five reported metrics.
The strongest zero-shot competitor is GPT-4o (F1~=~0.629, MCC~=~0.493),
followed by Claude Sonnet~4.5 (F1~=~0.598, MCC~=~0.446) and CodeLlama-13b
(F1~=~0.431, MCC~=~0.171).
Relative to GPT-4o, RetroBug improves Bug-Fix F1 by 14.7 points and MCC
by 0.205 points; relative to Claude, the gains are 17.8 and 0.252 points
respectively.

The performance gap is driven primarily by \emph{precision}.
All three zero-shot models exhibit a Bug-Fix bias: they achieve high recall
(0.625--0.800) but low precision (0.329--0.571), misclassifying a large
share of non-bug commits as bug-fixes.
Claude Sonnet~4.5, for example, labels 35 of 120 non-bug commits as
Bug-Fix (precision 0.478), while still recovering 32 of 40 true bug-fix
commits (recall 0.800).
CodeLlama-13b amplifies this pattern most severely: it predicts Bug-Fix
for 76 commits but only 25 are correct (precision 0.329), yielding accuracy
of just 0.585 despite moderate recall (0.625).
GPT-4o balances the trade-off best among zero-shot models (precision
0.571, recall 0.700), but still falls 16.2 precision points below
RetroBug (0.733).

RetroBug's advantage over single-pass prompting indicates that multi-step
backward CBR, knowledge-base retrieval, and structured diff analysis add
value beyond exposing the same evidence to a capable LLM in one shot.
The agent recovers 33 of 40 bug-fix commits (recall 0.825) while limiting
false positives to 12 of 120 non-bug commits---a precision--recall profile
that none of the zero-shot baselines match.
We evaluate LLM baselines on fold~0 only due to API cost and inference
latency; RetroBug results on the same fold are directly comparable because
they use the identical 160-commit test split.

% RQ2 — RetroBug vs CCS baselines (fold 0, per-class metrics)
% Requires: \usepackage{booktabs}
% Usage: % RQ2 — RetroBug vs CCS baselines (fold 0, per-class metrics)
% Requires: \usepackage{booktabs}
% Usage: % RQ2 — RetroBug vs CCS baselines (fold 0, per-class metrics)
% Requires: \usepackage{booktabs}
% Usage: \input{paper_tables/rq2_results.tex}

\begin{table}[t]
\centering
\scriptsize
\caption{RQ2: Per-class performance of RetroBug versus CCS zero-shot
LLM baselines on fold~0 ($n=160$).
NB~=~Non-Bug-Fix, BF~=~Bug-Fix.}
\label{tab:rq2}
\setlength{\tabcolsep}{3pt}
\renewcommand{\arraystretch}{0.9}
\begin{tabular}{@{}l ccc ccc cc@{}}
\toprule
 & \multicolumn{3}{c}{\textbf{Non-Bug-Fix (NB)}} & \multicolumn{3}{c}{\textbf{Bug-Fix (BF)}} & & \\
\cmidrule(lr){2-4} \cmidrule(lr){5-7}
\textbf{Method} & \textbf{P} & \textbf{R} & \textbf{F1} & \textbf{P} & \textbf{R} & \textbf{F1} & \textbf{Acc.} & \textbf{MCC} \\
\midrule
CCS (Claude)~\cite{zeng2025conventional}
  & 0.914 & 0.708 & 0.798 & 0.478 & 0.800 & 0.598 & 0.731 & 0.446 \\
\midrule
CCS (GPT-4o)~\cite{zeng2025conventional}
  & 0.892 & 0.825 & 0.857 & 0.571 & 0.700 & 0.629 & 0.794 & 0.493 \\
\midrule
CCS (CodeLlama-13b)~\cite{zeng2025conventional}
  & 0.819 & 0.571 & 0.673 & 0.329 & 0.625 & 0.431 & 0.585 & 0.171 \\
\midrule
\textbf{RetroBug (ours)}
  & \textbf{0.939} & \textbf{0.900} & \textbf{0.919} & \textbf{0.733} & \textbf{0.825} & \textbf{0.776} & \textbf{0.881} & \textbf{0.698} \\
\bottomrule
\end{tabular}

\vspace{0.3em}
{\footnotesize CodeLlama-13b: $n=159$ (one unparseable response excluded).}
\end{table}


\begin{table}[t]
\centering
\scriptsize
\caption{RQ2: Per-class performance of RetroBug versus CCS zero-shot
LLM baselines on fold~0 ($n=160$).
NB~=~Non-Bug-Fix, BF~=~Bug-Fix.}
\label{tab:rq2}
\setlength{\tabcolsep}{3pt}
\renewcommand{\arraystretch}{0.9}
\begin{tabular}{@{}l ccc ccc cc@{}}
\toprule
 & \multicolumn{3}{c}{\textbf{Non-Bug-Fix (NB)}} & \multicolumn{3}{c}{\textbf{Bug-Fix (BF)}} & & \\
\cmidrule(lr){2-4} \cmidrule(lr){5-7}
\textbf{Method} & \textbf{P} & \textbf{R} & \textbf{F1} & \textbf{P} & \textbf{R} & \textbf{F1} & \textbf{Acc.} & \textbf{MCC} \\
\midrule
CCS (Claude)~\cite{zeng2025conventional}
  & 0.914 & 0.708 & 0.798 & 0.478 & 0.800 & 0.598 & 0.731 & 0.446 \\
\midrule
CCS (GPT-4o)~\cite{zeng2025conventional}
  & 0.892 & 0.825 & 0.857 & 0.571 & 0.700 & 0.629 & 0.794 & 0.493 \\
\midrule
CCS (CodeLlama-13b)~\cite{zeng2025conventional}
  & 0.819 & 0.571 & 0.673 & 0.329 & 0.625 & 0.431 & 0.585 & 0.171 \\
\midrule
\textbf{RetroBug (ours)}
  & \textbf{0.939} & \textbf{0.900} & \textbf{0.919} & \textbf{0.733} & \textbf{0.825} & \textbf{0.776} & \textbf{0.881} & \textbf{0.698} \\
\bottomrule
\end{tabular}

\vspace{0.3em}
{\footnotesize CodeLlama-13b: $n=159$ (one unparseable response excluded).}
\end{table}


\begin{table}[t]
\centering
\scriptsize
\caption{RQ2: Per-class performance of RetroBug versus CCS zero-shot
LLM baselines on fold~0 ($n=160$).
NB~=~Non-Bug-Fix, BF~=~Bug-Fix.}
\label{tab:rq2}
\setlength{\tabcolsep}{3pt}
\renewcommand{\arraystretch}{0.9}
\begin{tabular}{@{}l ccc ccc cc@{}}
\toprule
 & \multicolumn{3}{c}{\textbf{Non-Bug-Fix (NB)}} & \multicolumn{3}{c}{\textbf{Bug-Fix (BF)}} & & \\
\cmidrule(lr){2-4} \cmidrule(lr){5-7}
\textbf{Method} & \textbf{P} & \textbf{R} & \textbf{F1} & \textbf{P} & \textbf{R} & \textbf{F1} & \textbf{Acc.} & \textbf{MCC} \\
\midrule
CCS (Claude)~\cite{zeng2025conventional}
  & 0.914 & 0.708 & 0.798 & 0.478 & 0.800 & 0.598 & 0.731 & 0.446 \\
\midrule
CCS (GPT-4o)~\cite{zeng2025conventional}
  & 0.892 & 0.825 & 0.857 & 0.571 & 0.700 & 0.629 & 0.794 & 0.493 \\
\midrule
CCS (CodeLlama-13b)~\cite{zeng2025conventional}
  & 0.819 & 0.571 & 0.673 & 0.329 & 0.625 & 0.431 & 0.585 & 0.171 \\
\midrule
\textbf{RetroBug (ours)}
  & \textbf{0.939} & \textbf{0.900} & \textbf{0.919} & \textbf{0.733} & \textbf{0.825} & \textbf{0.776} & \textbf{0.881} & \textbf{0.698} \\
\bottomrule
\end{tabular}

\vspace{0.3em}
{\footnotesize CodeLlama-13b: $n=159$ (one unparseable response excluded).}
\end{table}


\begin{tcolorbox}[colback=gray!8, colframe=gray!50,
  boxrule=0.5pt, arc=2pt, left=4pt, right=4pt,
  top=3pt, bottom=3pt]
\textbf{Answer to RQ2:}
On fold~0, RetroBug (Bug-Fix F1~=~0.776, MCC~=~0.698) outperforms the
best zero-shot LLM baseline, GPT-4o (F1~=~0.629, MCC~=~0.493), by 14.7
F1 points and 0.205 MCC points.
Zero-shot models achieve competitive recall but suffer from low Bug-Fix
precision (0.329--0.571), whereas RetroBug combines high recall (0.825)
with substantially higher precision (0.733).
These results show that agentic backward CBR and tool use improve
classification beyond single-pass LLM prompting on the same inputs.
\end{tcolorbox}
